% BOOK-FILE-ID: f-6dfda5da3083
% subsubsection: 極形式による図形的解釈

\subsubsection{複素平面で見る}
  前節では$\frac{a_n-\alpha}{a_n-\beta}$という一見突飛な置き換えで,一次分数型を等比型に帰着させた.
  この置き換えに図形的な意味はないのだろうか.
  特性方程式の解が虚数になる場合を例に,複素平面の上で漸化式を眺めてみよう.
  結論から述べれば,このとき一次分数型漸化式は「回転」として解釈できるのだ.

  \paragraphbf{特性方程式の解が虚数のとき}
  $p,q,r,s$が実数で,特性方程式
  \[
    rx^2+(s-p)x-q=0
  \]
  の2解が実数でないとする.
  実数係数の二次方程式なので,2解は互いに共役な複素数であり,\,$\beta=\overline{\alpha}$と書ける.
  前節と同じく
  \[
    b_n=\frac{a_n-\alpha}{a_n-\overline{\alpha}}
  \]
  と置くと,
  \[
    b_{n+1}=\lambda b_n,\qquad
    \lambda=\frac{p-r\alpha}{p-r\overline{\alpha}}
  \]
  となる.
  ここで$p,r$は実数だから$\overline{p-r\alpha}=p-r\overline{\alpha}$であり,\,$\lambda$の分子と分母は互いに共役である.
  共役な複素数の絶対値は等しいので,
  \[
    |\lambda|=1
  \]
  が分かる.
  したがって$\lambda$は極形式で$\lambda=\cos\theta+i\sin\theta$と書ける.

  $b_n=\rho_n(\cos\varphi_n+i\sin\varphi_n)$とすれば,複素数の掛け算は偏角の足し算だったから,
  \begin{align*}
    b_{n+1}
    &=(\cos\theta+i\sin\theta)\rho_n(\cos\varphi_n+i\sin\varphi_n)\\
    &=\rho_n\{\cos(\varphi_n+\theta)+i\sin(\varphi_n+\theta)\}
  \end{align*}
  である.つまり
  \[
    \rho_{n+1}=\rho_n,
    \qquad
    \varphi_{n+1}=\varphi_n+\theta
  \]
  となる.\footnote{偏角には$2\pi$の整数倍の違いがあるが,ここでは$\varphi_n=\varphi_1+(n-1)\theta$と選んでいる.}
  $b_n$は原点からの距離を変えずに,毎回$\theta$ずつ回転しているのだ.
  等比型への帰着は単なる計算の技巧ではなく,複素平面上の回転を表していたのである.

  \begin{bookpoint}
    \textbf{POINT}\par
    特性方程式の解が共役な虚数$\alpha,\overline{\alpha}$のとき,\,$b_n=\frac{a_n-\alpha}{a_n-\overline{\alpha}}$は絶対値1の複素数を掛けられ,原点のまわりを一定の角度ずつ回転する.
  \end{bookpoint}

  \paragraphbf{具体例}
  実際に
  \[
    a_1=\frac12,\qquad a_{n+1}=\frac{a_n+1}{1-a_n}
  \]
  で確かめてみよう.
  まずは実験である.順に計算すると
  \[
    \frac12\longrightarrow3\longrightarrow-2
    \longrightarrow-\frac13\longrightarrow\frac12\longrightarrow\cdots
  \]
  となり,4回で元に戻る.\footnote{どの項でも分母$1-a_n$は0でない.}
  この周期4の正体を調べよう.
  特性方程式は$x=\frac{x+1}{1-x}$,すなわち$x^2+1=0$なので,\,$\alpha=i,\ \beta=-i$である.
  そこで
  \[
    b_n=\frac{a_n-i}{a_n+i}
  \]
  と置くと,
  \[
    \lambda=\frac{1+i}{1-i}=\frac{(1+i)^2}{2}=i
  \]
  より$b_{n+1}=ib_n$を得る.
  $i=\cos\frac{\pi}{2}+i\sin\frac{\pi}{2}$だから,\,$b_n$は毎回反時計回りに$\frac{\pi}{2}$ずつ回転する.
  \[
    b_1\longrightarrow ib_1\longrightarrow-b_1
    \longrightarrow-ib_1\longrightarrow b_1\longrightarrow\cdots
  \]
  $\frac{\pi}{2}$の回転を4回繰り返せば一周するので,\,$a_n$も4回で元に戻るのだ.
  実数の上では不規則に見えた変化が,複素平面の上では一定角度の回転だったのである.

  \paragraphbf{小話：正接の加法定理}
  この漸化式に見覚えはないだろうか.
  正接の加法定理
  \[
    \tan\left(x+\frac{\pi}{4}\right)=\frac{\tan x+1}{1-\tan x}
  \]
  と見比べると,\,$a_n=\tan x_n$と置けば$x_{n+1}=x_n+\frac{\pi}{4}$となる.
  つまり,この漸化式は角度を$\frac{\pi}{4}$ずつ足す操作だったのだ.
  実際,\,$a_n=\tan x_n$を$b_n$に代入すると
  \[
    b_n=\frac{\sin x_n-i\cos x_n}{\sin x_n+i\cos x_n}
    =-(\cos2x_n+i\sin2x_n)
  \]
  となるので,\,$x_n$が$\frac{\pi}{4}$進むと$b_n$の偏角は$\frac{\pi}{2}$進む.
  一次分数型の回転と,三角関数の加法定理は同じものを見ていたのである.
  三角関数を使う解法は\sectionref{節：三角関数型}でも扱う.
